\section{Introduction}

Breast cancer is the most commonly diagnosed cancer among women worldwide, with over 2.3 million new cases reported annually \cite{WHO2025}. Early detection plays a crucial role in improving patient outcomes. Currently, X-ray mammography is the standard screening method; however, it has important limitations. Its sensitivity is significantly reduced in women with dense breast tissue, and the use of ionizing radiation makes it less suitable for repeated or early screening \cite{VonEuler-Chelpin2019, Bakker2019}. These limitations have motivated the development of alternative, non-invasive and radiation-free imaging approaches.

\subsection{Background}

Within this context, novel imaging approaches are being developed for radiation-free breast cancer detection. This project is conducted in collaboration with Early Warning Scan (EWS), a company that is developing a non-invasive screening technology based on optical and biomechanical measurements of breast deformation \cite{EarlyWarningScanWebsite}. The underlying idea is that mechanical differences between healthy and cancerous tissue can be detected through small variations in breast surface motion during controlled excitation. Malignant tumors are known to be significantly stiffer than surrounding tissue, which forms the basis for several biomechanical imaging techniques, such as elastography and Digital Image Elasto-Tomography (DIET) \cite{Samani2007, Patel2022, Peters2004}. In these methods, external mechanical excitation is applied to the breast, and the resulting deformation patterns are analyzed to infer internal tissue properties.

To support the development and validation of such imaging techniques, computational models of the breast are increasingly used. In particular, finite element method (FEM) models provide a powerful framework for simulating the mechanical behaviour of breast tissue under different loading conditions \cite{Samani2001, Hsu2011}. These models enable controlled investigation of deformation patterns and can be used to generate synthetic data for the development and evaluation of imaging algorithms.

Over time, FEM breast models have evolved from simplified homogeneous representations to more advanced heterogeneous and multi-component models. Early models often used simplified material assumptions or focused on limited anatomical components, which can restrict how well realistic tissue behaviour is represented \cite{Samani2001, Rajagopal2006, Ruggiero2014}. More recent work incorporates multiple anatomical components, such as adipose tissue, fibroglandular tissue, skin, and muscle, each with distinct mechanical properties \cite{Chen2024, Mazier2024}. In addition, image-based approaches using MRI or CT data have enabled the construction of more realistic, patient-specific geometries \cite{Hsu2011, Sturgeon2016}.

Despite these developments, important limitations remain. Many existing models still rely on simplified geometries and do not fully capture the anatomical variability and internal structural complexity of the breast. However, breast tissue is inherently heterogeneous and exhibits complex mechanical behaviour, with variations in material properties both between and within tissue types \cite{Ramiao2016, Goodbrake2022, Galbreath2025}. These simplifications can limit the accuracy of simulated deformation patterns and reduce the effectiveness of such models in applications such as imaging validation and AI training.

\subsection{Objectives}

The aim of this project is to improve the anatomical and biomechanical realism of an existing finite element breast model while keeping the workflow reproducible enough for systematic comparison. During the project, the main development branch shifted toward a COMSOL-based pipeline because this made it possible to automate constructive geometry, region selections, and post-processing for multiple model variants \cite{COMSOLMultiphysicsWebsite}. The final COMSOL workflow is organized as a self-contained source-case pipeline, so the active COMSOL route can be generated and checked without relying on the older FEBio-oriented package at runtime \cite{FEBioWebsite}.

The central research question of this study is:

\textit{How do controlled improvements in chestwall support, glandular composition, outer breast shape, Cooper ligament representation, and tumor/lesion inclusion affect simulated breast deformation patterns in a heterogeneous finite element breast model?}

To address this question, the project focuses on six connected model-development steps: establishing a mild dynamic baseline, introducing a volume-preserving curved chestwall representation, adding a realistic glandular reference, defining matched asymmetry and no-Cooper controls, preparing Cooper ligament support sensitivities, and adding an initial tumor/lesion screening route. The emphasis is on identifying which model components are already report-ready and which remain exploratory.
